\section{Induction as a search action}

\subsection{Generate a justified induction problem}

More ground reasoning is not a substitute for a structural argument. Forge
should inspect recursive definitions occurring in the goal, their equation
theorems, and the arguments on which recursive calls decrease. This yields a
small, ranked set of induction candidates. Structural induction on a list or
natural number is the first case; a function-specific induction theorem or a
well-founded recursor is used when the recursion requires it.

The algorithm must preserve the distinction between discovering an induction
hypothesis and assuming an arbitrary stronger assertion. A generalization
creates a \emph{new theorem to prove}, with its own base and step obligations.
Only the checked induction principle closes that theorem. An open induction
branch cannot contribute its conjectured invariant as a global fact.

A concrete procedure is:
\begin{enumerate}
\item Read the target and selected hypotheses as a dependency graph. Identify
recursive calls and the local variables in their changing argument positions.
\item Rank induction variables using occurrence in decreasing arguments,
recursor size, and expected branch count. Start with one variable, then bounded
pairs or a function-specific induction principle.
\item Form a generalization set from variables changed by the recursive calls.
Take its dependent-type closure, so reverting a variable does not strand local
declarations whose types depend on it.
\item Revert the selected declarations, apply the actual recursor or induction
theorem, introduce branch parameters and induction hypotheses, and unfold only
relevant equation lemmas.
\item Solve the resulting AND-branches cooperatively. On failure, retain the
residual goal to suggest a wider generalization or a missing auxiliary lemma,
but do not reuse the unproved assertion as evidence.
\end{enumerate}

A ``safe'' transformation is not necessarily cheap. Even semantically
appropriate induction can create an enormous proof search. Put a hard bound
on branch count and total candidate recursors before invoking it.

\subsection{Worked accumulator example}

Define, mathematically,
\begin{align*}
 \operatorname{revAcc}([],a)&=a,\\
 \operatorname{revAcc}(x::s,a)&=\operatorname{revAcc}(s,x::a).
\end{align*}
The natural target is
$\operatorname{revAcc}(s,[])=\operatorname{reverse}(s)$.
A weak induction hypothesis talks only about the accumulator $[]$, whereas the
recursive call uses $x::a$. It does not match.

The changed-argument analysis identifies the accumulator, suggesting
\begin{equation}\label{eq:rev-inv}
 \forall a,\quad \operatorname{revAcc}(s,a)
       =\operatorname{reverse}(s)\mathbin{+\!+}a.
\end{equation}
The empty-list case follows from the equations. In the cons case, the induction
hypothesis instantiated at $x::a$ gives
\[
 \operatorname{revAcc}(x::s,a)
 =\operatorname{reverse}(s)\mathbin{+\!+}(x::a).
\]
The recursive equation for reverse and associativity of append identify this
with $\operatorname{reverse}(x::s)\mathbin{+\!+}a$.
Specializing the now-proved generalized theorem at $a=[]$ gives the original
target. The complete proof idea is the quantified accumulator invariant, not
an unexplained instruction to ``try induction harder.''

The library may already contain a theorem that makes this example trivial.
That does not invalidate the generalization mechanism, but it changes the
benchmark. Tests of structural synthesis must use fresh definitions or control
which equivalent high-level lemmas are available. Never count retrieval of an
already-proved target as evidence that induction synthesis worked.

\subsection{Failure-guided auxiliary lemmas}

After an unsuccessful step, compare the recursive call with the induction
hypothesis. Mismatched accumulator arguments suggest generalization;
mismatched compositions suggest a fusion lemma; polynomial differences suggest
an invariant template. A candidate lemma is ranked by term size, number of new
quantifiers, and whether its recursive calls have a clear decreasing argument.
Try at most a bounded number at each size.

For a function $f$ with an accumulator update $a\mapsto T(a,x)$, templates can
include
\[
 f(s,a)=H(s,a),\qquad
 \mu(f(s,a))=\mu(a)+J(s),\qquad
 R(f(s,a),s,a).
\]
The template vocabulary is built from the goal, equation theorems, and a small
approved operator set. The resulting lemma is an OR-alternative with a separate
AND proof tree. If its proof closes, inject its specialization into the parent.
Otherwise discard it or widen the budget. This avoids mutually recursive
``helper lemmas'' becoming a circular proof.

\section{Polynomial invariant synthesis with exact validation}

\subsection{From recurrence data to a theorem candidate}

For a state transition $T$ and starting state $s_0$, an equality invariant
$I(s)=0$ is justified by
\[
 I(s_0)=0,\qquad I(T(s))=I(s).
\]
The second identity is a convenient sufficient condition. More generally it
would suffice to prove $I(s)=0\Rightarrow I(T(s))=0$, but that is not the
certificate language implemented here.

The prototype considers additive polynomial recurrences
\begin{equation}\label{eq:additive}
 n'=n+1,\qquad s'=s+p(n),\qquad (n_0,s_0)=(0,c),
\end{equation}
where $p\in\Q[n]$. It proposes a polynomial $Q(n)$ and sets
$I(n,s)=s-Q(n)$. The certificate obligations become
\begin{equation}\label{eq:fdiff}
 Q(0)=c,\qquad Q(n+1)-Q(n)=p(n).
\end{equation}
Given a degree bound $d$, sampling the recurrence at $d+1$ points and
interpolating suggests $Q$. \emph{Sampling is only discovery.} The validator
substitutes the polynomial transition into $I$ and compares exact coefficient
maps for the universal identity.

A production synthesizer should also support direct coefficient solving.
Write $Q(n)=\sum_{i=0}^{d}c_i n^i$, expand \eqref{eq:fdiff}, and solve the
resulting rational linear system. This removes numerical concerns and avoids
a dependence on the availability of inexpensive evaluation samples. The
interpolation implementation was chosen here because it makes the separation
between plausible evidence and universal checking particularly explicit.

\subsection{Example: sum of squares}

For $p(n)=(n+1)^2$ and $c=0$, the discovered expression is
\[
 Q(n)=\frac{n(n+1)(2n+1)}6.
\]
Equivalently take the denominator-cleared invariant
\[
 J(n,s)=6s-n(n+1)(2n+1).
\]
Its step difference is
\begin{align*}
 J(n+1,s+(n+1)^2)-J(n,s)
 &=6(n+1)^2\\
 &\quad-(n+1)(n+2)(2n+3)+n(n+1)(2n+1)\\
 &=0.
\end{align*}
Its base value is zero. Ordinary induction on the number of iterations
therefore proves the equality for every reachable state. No extrapolation from
finite data is involved after the two identities have been checked.

The tests include a hostile interpolation example: an added high-degree term
vanishes at all sampled points but violates the recurrence identity elsewhere.
The exact validator rejects the proposed invariant. This is an essential test
for any counterexample-guided or sample-based synthesis component.

\begin{proposition}[Invariant certificate soundness]
Let $s_{k+1}=T(s_k)$. If $I(s_0)=0$ and $I\circ T=I$, then
$I(s_k)=0$ for every $k\in\N$.
\end{proposition}
\begin{proof}
The claim for $k=0$ is the base condition. If $I(s_k)=0$, then
$I(s_{k+1})=I(T(s_k))=I(s_k)=0$. Apply induction.
\end{proof}

\subsection{Generalizing the invariant engine}

For a polynomial state vector $\vec x$, use the template
\[
 I(\vec x)=\sum_{|\alpha|\leq d}c_\alpha\vec x^\alpha.
\]
Because $T$ is fixed, the coefficients of $I(T(\vec x))-I(\vec x)$ are linear
in the unknown $c_\alpha$. A rational nullspace computation discovers conserved
polynomials, with the additional linear constraint $I(s_0)=0$. Reject the
identically zero invariant and rank remaining candidates by relevance to the
postcondition. This is an immediately implementable extension, although the
shipped discovery routine handles only \eqref{eq:additive}.

A more expressive certificate permits a vector of invariant equalities:
\[
 I_i(T(\vec x))=\sum_j h_{ij}(\vec x)I_j(\vec x).
\]
This proves preservation of their common zero set. Bounded-degree multiplier
search reduces the check to exact identities, but jointly discovering the
$I_j$ and multipliers is no longer the same easy linear problem. Separate
candidate generation from multiplier certification, or fix one side at a time;
do not describe the joint bilinear search as linear algebra.

Inequality invariants require base and preservation \emph{inequalities}. The
PolyCone plugin can certify these under transition guards and existing
invariants. Termination remains a different obligation: an invariant does not
prove that a loop or recursive definition terminates. For already accepted
Lean definitions, use their existing recursion principles; for a proposed new
recursive program, termination must be proved separately.

\section{Existential witnesses and counterexample-guided synthesis}

\subsection{Construct witnesses in the correct scope}

For a goal $\exists y:\tau,\ P(\vec x,y)$, the useful output is a term $t$ of
\emph{exactly} type $\tau$ and a proof of $P(\vec x,t)$. Search begins with
terms already in scope, constructor applications, affine expressions suggested
by equations, and terms built from the relevant theorem signatures. Candidate
terms must not depend on variables introduced inside a later branch or a
future universal quantifier.

For
\[
 \forall n\in\N,\quad\exists m\in\N,\quad m>n\ \wedge\ m\equiv0\pmod2,
\]
a bounded affine grammar can propose $m=2(n+1)$. The remaining proof obligations
are natural-number arithmetic and divisibility. Construct the existential
proof explicitly after those obligations close. A real-valued LP solution for
$m$ would not be a natural-number witness; domain checks are part of the
construction, not an optional verification pass.

\subsection{A concrete CEGIS loop}

Counterexample-guided inductive synthesis (CEGIS) alternates candidate proposals
and attempted refutations of their universal correctness. For an affine
integer grammar $t(\vec x)=c_0+\sum_i c_i x_i$, enumerate bounded integer
coefficient vectors or ask an integer solver for coefficients that satisfy a
finite sample set. Once a candidate is found, ask the cooperative prover to
establish $\forall\vec x,\ P(\vec x,t(\vec x))$ under the actual assumptions.

\begin{Code}
synthesize(grammar, samples, goal, budget):
  repeat within term-size and coefficient bounds:
    t := candidate fitting every validated sample
    result := prove_universally(goal specialized at t)
    if result is a checked proof:
      return explicit witness t and proof
    if result contains a validated counterexample:
      add it to samples
    otherwise:
      schedule a different candidate or a larger proof budget
  return unknown
\end{Code}

An external model is not automatically a Lean counterexample. Recheck its
values against the encoded operations and assumptions, especially casts,
truncating subtraction, division, and finite types. A spurious model should
refine the encoding or be ignored; it must not eliminate a valid candidate.
Likewise, failure to prove a candidate does not mean it is false.

This module is specified, not implemented in the package. A reasonable first
implementation restricts the grammar to constructors, affine arithmetic, and
finite tuples, reusing \code{omega}, \code{linarith}, and \grind\ for residual
proofs. The current reference explicitly does not describe \code{omega} as a
complete procedure for every integer-linear problem \cite{lean-tactics}; the
integration should preserve ordinary tactic failure semantics.

\subsection{Models guide search; they do not certify theorems}

The same discipline applies to learned premise selection, algebraic pattern
recognition, and suggestions from a computer algebra system. A proposed
factorization becomes useful only after the exact identity is proved. A
proposed induction invariant becomes useful only after base and preservation
are proved. A proposed witness becomes useful only after its universally
scoped specification is proved. This uniform boundary makes the architecture
compatible with more ambitious search methods without enlarging the logical
trust base merely because a search component is complicated.
